%%%%%%%%%%%%%%%%%%%%%%%%%%%%%%%%%%%%%%%%%%%%%%%%%%
% Metadata
%%%%%%%%%%%%%%%%%%%%%%%%%%%%%%%%%%%%%%%%%%%%%%%%%%
% id: jmo-2026-y-2
% title: 2026年 JMO 予選 第2問
% tags: [整数]
% difficulty: 5
% source: https://www.imojp.org/archive/mo2026/jmo/problems/jmo36yqa.pdf

%%%%%%%%%%%%%%%%%%%%%%%%%%%%%%%%%%%%%%%%%%%%%%%%%%
% Preamble
%%%%%%%%%%%%%%%%%%%%%%%%%%%%%%%%%%%%%%%%%%%%%%%%%%
\documentclass[fleqn]{ltjsarticle}

\usepackage{common}
\loadcommonpreamble

% ヘッダー
\lhead{\textbf{2026年 JMO 予選 第2問\footnote[1]}}
\rhead{\repeatstars{5}}

%%%%%%%%%%%%%%%%%%%%%%%%%%%%%%%%%%%%%%%%%%%%%%%%%%
% Document
%%%%%%%%%%%%%%%%%%%%%%%%%%%%%%%%%%%%%%%%%%%%%%%%%%
\begin{document}

\footnotetext[1]{Copyright \textcopyright\space 2025 by Mathematical Olympiad Foundation of Japan.  \\ 著作権は数学オリンピック財団に帰属します．また，問題は以下から引用しています．\\ \url{https://www.imojp.org/archive/mo2026/jmo/problems/jmo36yqa.pdf}}

\begin{problembox}{問題}
  正の整数 $n$ であって，$\gauss{\sqrt{2026}n}$ が $n$ で割り切れないようなもののうち，最小のものを求めよ．
  ただし，実数 $r$ に対して $r$ 以下の最大の整数を $\gauss{r}$ で表す．たとえば，$\gauss{3.14}=3, \gauss{5}=5$ である．

  \syutten{数学オリンピック予選}
\end{problembox}

\noindent
\kaib\quad
$45<\sqrt{2026}<46$ より，$0<\alpha<1$ の実数 $\alpha$ を用いて
\begin{align*}
  \sqrt{2026}=45+\alpha
\end{align*}
と表せる．このとき
\begin{align*}
  \gauss{\sqrt{2026}n}=\gauss{45n+\alpha n}=45n+[\alpha n]
\end{align*}
と書ける．
\begin{itemize}
  \item [] $0<\alpha n<1$ のとき \\
    $45n+\gauss{\alpha n}=45n$ より，
    $\gauss{\sqrt{2026}n}$ は $n$ で割り切れる．
  \item [] $1<\alpha n$ のとき \\
    $\alpha n<n$ であるから，$\gauss{\alpha n}$ は $n$ で割り切れない． \\
    ゆえに，$\gauss{\sqrt{2026}n}$ も $n$ で割り切れない．
\end{itemize}
以上から，$1<\alpha n$ をみたす最小の $n$ が求める答えである． \\
$1<\alpha n$ を解くと
\begin{align*}
  &(\sqrt{2026}-45)n>1 \Y n>\sqrt{2026}+45
\end{align*}
$90<\sqrt{2026}+45<91$ より，求める値は $\boldsymbol{n=91}\kotae$

% \noindent
% \sankoua\quad

\end{document}